% Methods. Describes the pipeline in bench/ (PLAN.md has the operational detail). Keep in step with bench/config.yaml.
\section{Supplementary methods}
\label{sec:methods}

\subsection{Datasets}
The \nm{nDatasets} public datasets (Table~\ref{tab:datasets}; accessions, licences and conversion notes in
\texttt{datasets.yaml}) are each downloaded from the primary repository at a fixed accession and file, checked against a
recorded SHA-256, and converted to a canonical CSC matrix: features as rows, cells as columns, unsigned 32-bit counts, no
explicit zeros, and the source's feature identifiers and barcodes. Non-integer matrices are rejected, not rounded. Each
count matrix of a multimodal dataset is a separate unit. Bulk RNA-seq gene counts and a bag-of-words corpus are run as an
extension.

\subsection{Formats}
Each format of Table~\ref{tab:formats} is written from the canonical matrix by its own library with default or
documented recommended options and the compression listed there; R \texttt{dgCMatrix} is saved with \texttt{saveRDS},
loom only where its dense chunks fit the run budget, and the CSC arrays raw and with zstd level 19. The generic-codec
control, Zarr v3 (zarr-python 3.1), stores the CSR \texttt{data} and \texttt{indices} in inner chunks of 49,152 elements
(a \sone{} chunk) in shards of 64, each compressed with Blosc (zstd level 5, bit shuffling) in the narrowest unsigned type
that holds it, with \texttt{indptr} as one chunk and names as gzipped text. Every file is read back and compared with the
canonical matrix, names included; a failure is reported, not dropped. Above $2^{31}-1$ non-zeros, the limit of a
\texttt{dgCMatrix}, the RDS formats and BPCells are not applicable, and BPCells leaves the minibatch panel. Size is the
bytes of all the format's files per non-zero. \sone{} uses the \texttt{singlet} writer's defaults; the shared-pack variant
stores gene identifiers in a vocabulary file whose size is not counted.

\subsection{Read speed}
Each configuration runs in a fresh process pinned to $p$ cores: every format at $p=1$ and $p=\nm{threadsHeadline}$
(headline), and \sone{}, gzip AnnData and the raw CSC arrays also at $p\in\{1,2,4,8\}$; three timed repeats, median
reported, every time kept. An untimed write and read of a 60-by-40 matrix in the same format first moves imports and
first-use set-up out of the timing (R formats load BPCells first); the first cold repeat then takes a median of
\nm{prot-first-over-median} times the median of three (\nm{prot-n} configurations). Cold reads evict the file from the
page cache (\texttt{posix\_fadvise(POSIX\_FADV\_DONTNEED)}) before every repeat; the disk controller's cache cannot be
flushed from a job (a just-written file reads at about 5\,GB/s after 16\,GB of other I/O, against 1.4\,GB/s from disk)
and is in the same state for every repeat and format. Warm reads, at $p=\nm{threadsHeadline}$ only, follow one untimed
read. In Python a read longer than 60\,s is timed once; in R every repeat is timed. Matrices above $4\times10^8$
non-zeros (\texttt{decode.big\_nnz}) are timed cold and warm at $p=\nm{threadsHeadline}$ from local disk only, for the
CSC arrays, Cell Ranger HDF5, AnnData uncompressed and gzip, AnnData Zarr, TileDB-SOMA and \sone{} (also with the shared
pack).

The timed interval runs from opening the file to the library's object (a SciPy matrix with names, an AnnData object, or
the R equivalent); the conversion and check against the canonical matrix, done at the write stage, are not timed.
\sone{} is timed in its stored gene order, the reader's default, with gene names permuted to match, and in source gene
order. The \sone{} codec alone is timed from a member in memory to CSC arrays, through the Python binding and a C++
program, for each instruction set (scalar, AVX2, AVX-512) and both gene orders, recording which kernel ran (an
unsupported request falls back to the next).

\subsection{Minibatch loading}
The loader panel has four datasets, from small to the largest tier, from droplet, combinatorial and plate protocols. A
workload is a seeded sequence of 512-cell minibatches, each eight blocks of 64 consecutive cells at random offsets, as
block-shuffling loaders read~\cite{dascenzo2025scdataset}. Loaders: AnnData in memory and backed, AnnData Zarr read
lazily, scDataset over backed AnnData, TileDB-SOMA ML, BPCells (R, timed to an R sparse matrix), and the \sone{} reader
(\texttt{singlet.torch.S1Stream}) in stored and source gene order; Python loaders return a
\texttt{torch.sparse\_csr\_tensor}. scDataset and TileDB-SOMA ML use their own samplers set to the same workload (blocks
of 64, one minibatch per read, no larger shuffle buffer). The \sone{} reader plans a minibatch's byte ranges from the chunk
index and sends up to eight at a time over HTTP (\texttt{urllib3}; \texttt{singlet} 1.3.2 default, 1.3.3 sends 32), or
one at a time; decodes all of its runs in parallel into one CSR in one C++ call outside Python's interpreter lock; and,
as \texttt{S1Cells} does by default, sends the next minibatch's requests while one decodes. TileDB-SOMA ML prefetches with
its own reader; the others read one minibatch after another.

Files are served from the node's disk, from the cluster's shared NFS filesystem, and by a local S3-compatible store
(\texttt{rclone serve s3}) at 0.1, 1 or 10\,Gbit/s with 1\,ms latency, or 1\,Gbit/s with 30\,ms. Bandwidth is one token
bucket over all connections (our proxy, \texttt{bwproxy.py}) holding at most 5\,ms of the tier's rate or 64\,KiB,
whichever is larger, so concurrent readers share the tier rather than multiply it; Toxiproxy adds the latency to every
response. At each tier the rate without latency is measured with a 1\,GiB transfer over one and over $p$ connections: a
median of \nm{delivered-s3_0.1g}, \nm{delivered-s3_1g} and \nm{delivered-s3_10g}\,Gbit/s over $p$ at the three tiers. A
tier not reached over $p$ connections (\nm{cappedTiers}) is flagged. Loaders that cannot read by range download the object
first, and that time counts; above the M tier they are not run over the network tiers, and the object's transfer time at
the tier's bandwidth is reported instead (marked in Table~\ref{tab:minibatch}).

From a new process with a cold cache and imports done, we record the time to open, to read the first minibatch and to
build its tensor (three seeds; medians), then, for one seed, the sustained rate over 200 minibatches after 10 warm-up
ones, and peak resident memory. Below $211\times512$ cells the sustained workload reads some cells twice, so ratios of
sustained rates use only larger datasets. Link share is the sustained rate times the file's bytes per cell over the
bandwidth. From local disk each loader reads its own file, evicted first; the RAID cache cannot be cleared, so of the
three \sone{} runs on one file (stored order, source order, one request at a time) only the first, which is reported, is
cold.

\subsection{Random-access costs, latency and a real object store}
\emph{Access costs} are computed, not timed, for 200 minibatches per matrix: for \sone{} from the chunk index (the ranges
\texttt{S1Stream} plans, adjacent ranges merged), for the HDF5 formats from the chunk table of \texttt{data},
\texttt{indices} and \texttt{indptr} (\texttt{h5py} \texttt{get\_chunk\_info}), and for AnnData Zarr from its chunk
object sizes; a request is one contiguous range or one object. Over-read is bytes fetched divided by the minibatch's
non-zeros times the file's bytes per non-zero. For the chunk-size sweep the \sone{} member of three matrices is
re-encoded with chunk budgets of 4,096 to 393,216 non-zeros.

\emph{Latency and requests in flight.} The emulated store at 1\,Gbit/s adds 0, 5, 10, 30, 100 or 200\,ms to every
response and counts the bytes and HTTP requests it serves. Each point is one fresh process: the first minibatch, then
60 minibatches after 5 warm-up ones, each timed. The in-flight sweep runs the \sone{} reader at 30 and 100\,ms with 1
to 64 requests at once and 0 to 8 minibatches sent ahead. The DataLoader comparison runs \texttt{S1Cells}
(\texttt{singlet} 1.3.3 defaults) and scDataset with 0, 2 and 4 workers (equal shares of the cores,
\texttt{prefetch\_factor} 2), from local disk and from the store at 1\,ms; scDataset also runs with its default
\texttt{fetch\_factor} of 16, which reads 16 minibatches at once and shuffles them in memory.

\emph{A real object store.} The files were read from a Cloudflare R2 bucket over the public internet from the timing
node, through a Cloudflare Worker serving the bucket by range (\texttt{GET} with \texttt{Range}, \texttt{HEAD}; source
in \texttt{bench/r2worker}), as an unsigned HTTP reader reads a public bucket. The link was measured with 20 one-byte
range requests to an object and 20 to a route the Worker answers without the bucket, and by reading the \sone{} file on
one connection and in eight ranges at once. Each loader read three seeds as above; the \sone{} reader also with 32 and 64
requests at once and 4 and 8 minibatches ahead. An epoch's cost is its minibatches (cells / 512) times the requests and
bytes per minibatch counted at the emulated store without latency, at R2's prices (Class B reads \$0.36 per million, no
egress) and S3 Standard's in us-east-1 (\texttt{GET} \$0.004 per 10,000, transfer to the internet \$0.09 per GB) on 30
October 2026, without free tiers; downloading the h5ad is one request and the file's size.

\emph{Writing} is timed in a fresh process per format, library primed as for reads, from the canonical matrix in memory to
local disk; memory is the peak resident memory of the process and its children minus that before writing.

\subsection{Payload ladder}
Payload size, in bits per non-zero over row indices and values only, of: 32-bit CSC; LEB128 varints; row gaps and values
as LEB128 in separate streams with zstd at levels 3 and 19; and the \sone{} chunk payload (rANS tables, row keys and
chunks, without header and names).

\subsection{Decoding limits}
\label{sec:mkernels}
The \texttt{kernels} stage (\texttt{bench/cpp/kernel\_bench.cpp}, \texttt{-O3 -march=x86-64-v3}) reads each canonical
matrix up to $4\times10^8$ non-zeros (\nm{kern-n} matrices) and expresses each cell's row indices as steps (row $-$
previous row $-$ 1, as \sone{} does). Each code stores two streams, steps and values: raw 32-bit (decoding is a copy);
fixed-width bytes (the narrowest whole number of bytes holding every value); LEB128; Stream VByte (two-bit length codes,
one byte shuffle per four values); vertical bit-packing in blocks of 256 values, eight 32-bit lanes and one width per
block, unpacked with AVX2 shifts; order-0 rANS with one static 12-bit frequency table per stream, four interleaved states,
symbols $\min(x,255)$ and an LEB128 escape stream; and zstd at levels 3 and 19 on the LEB128 streams. Decoding returns
32-bit row indices and values, so step codes include the pass back to row indices in the timed work. Each code's first
decode is compared byte for byte with the input. The \sone{} member is decoded on one thread with the AVX2 kernel in
source order; its size is the whole member, index and tables included. The \texttt{s1order} stage, run later on the same
node, decoded each member again in source and in stored order (row ids in stored gene space plus the gene permutation, no
per-cell sort) and checked that the stored-order output, mapped and sorted, equals the source-order output. Its
source-order rates were \nm{s1order-rerun-pct}\% of the first run's (median), so the stored-order rate used is the first
run's source-order rate times the ratio measured within the second run. \texttt{store\_bound} only fills the two output
arrays (\texttt{memset}). Every code runs pinned to one core, seven timed repeats after two warm-ups, median; all
matrices ran in sequence in one job on one node (\nm{cpu}).

For Fig.~\ref{M-fig:floor}c, a code with $b$ bits per non-zero and one-core rate $D_1$, decoding on $p$ cores from a link of
bandwidth $B$, delivers $R = (b/B + 1/(pD_1))^{-1}$ (Eq.~\ref{M-eq:reff}), so $R/p$ depends only on $x = B/p$. The ideal
code has the best count model's code length and the \texttt{store\_bound} rate of the same matrix; the practical optimum
has \sone's size and the rate of the fastest entropy-coded stream (rANS or zstd) on that matrix. Curves are medians over
matrices of the per-matrix ratio to the ideal code.

\subsection{Count models and the floor}
The code length of a model is $\sum_{g,c} -\log_2 p(x_{gc})$ over all entries, zeros included, per non-zero, under
two-fold cross-validation over a random split of cells (fit on one half, score the other, then swap). Describing the model
costs nothing, so each number is optimistic for a coder that would have to ship it. Models: \emph{Poisson},
$x_{gc}\sim\mathrm{Poisson}(s_c\mu_g)$ with $s_c$ the cell's total over the mean total; \emph{negative binomial},
$x_{gc}\sim\mathrm{NB}(s_c\mu_g,\theta_g)$ with $\theta_g$ by moments, smoothed over mean expression as in
sctransform~\cite{hafemeister2019normalization}; \emph{gene $\times$ depth}, per gene and per quantile bin of cell depth
(32 bins) a presence probability and (in 4 coarser bins) a distribution over values 1--63 with an Elias-gamma escape
above, Krichevsky--Trofimov smoothed~\cite{krichevsky1981performance}; and \emph{gene $\times$ depth $\times$ cluster},
the same also conditioned on the cell's cluster. Clusters come from Leiden~\cite{traag2019louvain} on the training half
(50 principal components of log-normalised counts; at most 200{,}000 training cells, the rest assigned to the nearest
centroid); test cells take the nearest training centroid and are charged the entropy of their label. Sparse cluster
contexts back off to gene $\times$ depth with a pseudo-count of 8. Zero terms of the parametric models are summed per
gene and depth bin at the bin's mean depth; on matrices of at most $5\times10^7$ entries the exact sum differs by at most
\nm{nbBinningErrPctMax}\%. A matrix's floor is the smallest of the four code lengths.

\subsection{Sampling estimate}
For 2{,}000 random test cells the profile is the cell's normalised counts shrunk toward its cluster,
$\pi_c = (x_c + \alpha\,\bar\pi_{k(c)})/(N_c + \alpha)$, with cluster profiles from the training half and $\alpha$ chosen
by binomial thinning~\cite{batson2019molecular,neufeld2024inference}: half of each cell's molecules are held out and their
likelihood under the shrunk profile of the other half is maximised. The entropy of $\mathrm{Multinomial}(N_c,\pi_c)$, with
$N_c$ the cell's total, is $-\log_2 N_c! + \sum_g \mathbb{E}[\log_2 x_{gc}!] - N_c \sum_g \pi_{gc}\log_2\pi_{gc}$ under
$x_{gc}\sim\mathrm{Binomial}(N_c,\pi_{gc})$, with the expectation summed exactly where $N_c\pi_{gc}<30$ and by a
second-order expansion above. Because $\pi_c$ uses the cell's own counts this is not a bound; it is reported for UMI
matrices only, per non-zero and as a fraction of the whole matrix's \sone{} payload per non-zero.

\subsection{Second-pass compression}
Each \sone{} file was compressed again with zstd \texttt{-{}-ultra -22}, xz \texttt{-9e}, brotli \texttt{-q 11} and
bzip2 \texttt{-9}; we report output over input size.

\subsection{Ecosystem}
For two datasets we timed, from a cold file, the calls of Fig.~\ref{M-fig:eco} and the usual readers of the same data:
\texttt{anndata.read\_h5ad}, \texttt{scanpy.read\_10x\_h5}, Seurat \texttt{Read10X\_h5} (a bare matrix) and
\texttt{CreateSeuratObject} over it (the object \texttt{singlet::load(as = "seurat")} returns), and zellkonverter
\texttt{readH5AD}. Python results are compared with the canonical matrix, R results with the matrix \texttt{Read10X\_h5}
returns, itself verified at the write stage. The C++ row times opening and decoding a member already in memory.

\subsection{Hardware, software and reproducibility}
Every timed job had 8 cores of a shared cluster node of one type (\nm{cpu}, \nm{cores} cores, \nm{memGB}\,GB memory,
node-local \nm{localDisk}), at most one of our timed jobs per node, other users' jobs beside it. Readers used
\nm{threadsHeadline} threads unless stated; versions are in Table~\ref{tab:formats}. The pipeline is \texttt{python -m
s1bench} in \texttt{papers/manuscripts/s1-format/bench} of the \texttt{singlet} repository (stages \texttt{fetch},
\texttt{canon}, \texttt{write}, \texttt{decode}, \texttt{ladder}, \texttt{floor}, \texttt{second}, \texttt{minibatch},
\texttt{eco}, \texttt{report}), with settings in \texttt{config.yaml}, the environment in \texttt{environment.yml}, the
versions that ran (conda export, BPCells and Toxiproxy releases, host records) in \texttt{results/}, and Slurm scripts in
\texttt{bench/slurm/clipper}. The \texttt{report} stage writes every plotted value to \texttt{generated/*.tsv} and every
number in the text to \texttt{generated/numbers.tex}; no result is typed by hand.
